\documentclass[../../main.tex]{subfiles}% 注意这里的文件路径不能用 ./main.tex ,否则用latexmk编译子文件会报错
\graphicspath{{\subfix{./image/}}} % 指定图片目录,后续可以直接使用图片文件名
% 注意这里的文件路径不能用 ../../image/ ,否则用latexmk编译子文件会报错

% 例如:
% \begin{figure}[H]
% \centering
% \includegraphics[scale=0.3]{图.png}
% \caption{}
% \label{figure:图}
% \end{figure}
% 注意:上述\label{}一定要放在\caption{}之后,否则引用图片序号会只会显示??.

\begin{document}

\section{$L^p$空间的范数公式}

\begin{theorem}
若 $f\in L^p(E)$（$1\leqslant p<+\infty$），则存在 $g\in L^{p'}(E)$ 且 $\|g\|_{p'}=1$，使得
\begin{align}
\|f\|_p = \int_E f(x)g(x)\mathrm{d}x. \label{eq::8u238rURU32ur8-9}
\end{align}
\end{theorem}
\begin{proof}
不妨设 $\|f\|_p \neq 0$.
\begin{enumerate}[(1)]
\item 当 $p=1$ 时，令 $g(x)=\operatorname{sign} f(x)$，则 $\|g\|_\infty=1$，且
\begin{align*}
\int_E f(x)g(x)\mathrm{d}x = \int_E |f(x)|\mathrm{d}x = \|f\|_1.
\end{align*}

\item 当 $1<p<+\infty$ 时，令
\begin{align*}
g(x) = \left( \frac{|f(x)|}{\|f\|_p} \right)^{p-1} \operatorname{sign} f(x),
\end{align*}
则
\begin{align*}
\int_E |g(x)|^{p'} \mathrm{d}x = \frac{1}{\|f\|_p^p} \int_E |f(x)|^p \mathrm{d}x = 1,
\end{align*}
且
\begin{align*}
\int_E f(x)g(x)\mathrm{d}x = \int_E |f(x)| \left( \frac{|f(x)|}{\|f\|_p} \right)^{p-1} \mathrm{d}x = \|f\|_p.
\end{align*}
\end{enumerate}

\end{proof}

\begin{theorem}
若 $f\in L^\infty(E)$，则
\begin{align*}
\|f\|_\infty = \sup_{\|g\|_1=1} \left\{ \left| \int_E f(x)g(x)\mathrm{d}x \right| \right\}.
\end{align*}
\end{theorem}

\begin{proof}
令 $\|f\|_\infty = M>0$，只需证明对任意 $\varepsilon>0$，存在 $g\in L^1(E)$ 且 $\|g\|_1=1$，使得
\begin{align*}
\int_E f(x)g(x)\mathrm{d}x > M - \varepsilon.
\end{align*}
由 $\|f\|_\infty=M$，存在 $E$ 的子集 $A$，$m(A)=a>0$，使得
\begin{align*}
|f(x)| > M - \varepsilon, \quad x\in A.
\end{align*}
令
\begin{align*}
g(x) = \frac{1}{a} \chi_A(x) \operatorname{sign} f(x),
\end{align*}
则
\begin{align*}
\|g\|_1 = \int_E |g(x)|\mathrm{d}x = \frac{1}{a} \int_E \chi_A(x)\mathrm{d}x = 1,
\end{align*}
且
\begin{align*}
\int_E f(x)g(x)\mathrm{d}x &= \frac{1}{a} \int_E f(x)\chi_A(x)\operatorname{sign} f(x)\mathrm{d}x \\
&= \frac{1}{a} \int_A |f(x)|\mathrm{d}x > M - \varepsilon.
\end{align*}

\end{proof}

\begin{remark}
若 $E=[0,1]$，$f(x)=x$，则 $\|f\|_\infty=1$。此时对任意 $g\in L^1(E)$ 且 $\|g\|_1=1$，有
\begin{align*}
\left| \int_0^1 f(x)g(x)\mathrm{d}x \right| \leqslant \int_0^1 x|g(x)|\mathrm{d}x < \int_0^1 |g(x)|\mathrm{d}x = 1,
\end{align*}
故当 $p=+\infty$ 时，不一定存在 $g\in L^1(E)$ 且 $\|g\|_1=1$，使得 $\|f\|_\infty = \int_E f(x)g(x)\mathrm{d}x$.
\end{remark}

\begin{theorem}\label{theorem:实变函数--定理6.19}
设 $g(x)$ 是 $E\subseteq\mathbb{R}^n$ 上的可测函数。若存在 $M>0$，使得对一切在 $E$ 上可积的简单函数 $\varphi(x)$，都有
\begin{align*}
\left| \int_E g(x)\varphi(x)\mathrm{d}x \right| \leqslant M \|\varphi\|_p,
\end{align*}
则 $g\in L^{p'}(E)$（$p'$ 是 $p$ 的共轭指标），且 $\|g\|_{p'}\leqslant M$.
\end{theorem}
\begin{proof}
\begin{enumerate}[(1)]
\item 当 $p>1$ 时，对 $|g(x)|^{p'}$，取具有紧支集的非负可测简单函数渐升列 $\{\varphi_k(x)\}$，使得
\begin{align*}
\lim_{k\to\infty} \varphi_k(x) = |g(x)|^{p'}, \quad x\in E.
\end{align*}
令 $\psi_k(x)=(\varphi_k(x))^{\frac{1}{p}} \operatorname{sign} g(x)$，则
\begin{align*}
\|\psi_k\|_p = \left( \int_E \varphi_k(x)\mathrm{d}x \right)^{\frac{1}{p}}.
\end{align*}
注意到
\begin{align*}
0 \leqslant \varphi_k(x) = (\varphi_k(x))^{\frac{1}{p}} (\varphi_k(x))^{\frac{1}{p}'} \leqslant (\varphi_k(x))^{\frac{1}{p}} |g(x)| = \psi_k(x)g(x),
\end{align*}
由假设可得
\begin{align*}
\int_E \varphi_k(x)\mathrm{d}x \leqslant \int_E \psi_k(x)g(x)\mathrm{d}x \leqslant M \|\psi_k\|_p,
\end{align*}
故
\begin{align*}
\int_E \varphi_k(x)\mathrm{d}x \leqslant M^{p'}.
\end{align*}
令 $k\to\infty$，得
\begin{align*}
\int_E |g(x)|^{p'} \mathrm{d}x \leqslant M^{p'}.
\end{align*}
\item 当 $p=1$ 时，不妨设 $g(x)\geqslant 0$。若 $g\notin L^\infty(E)$，则存在 $E$ 中的可测集列 $\{A_k\}$，$0<m(A_k)<+\infty$（$k=1,2,\cdots$），使得
\begin{align*}
g(x) \geqslant k, \quad x\in A_k, \quad k=1,2,\cdots.
\end{align*}
令 $\varphi_k(x)=\chi_{A_k}(x)$，则
\begin{align*}
\frac{\int_E \varphi_k(x)g(x)\mathrm{d}x}{\|\varphi_k\|_1} \geqslant \frac{k m(A_k)}{m(A_k)} = k, \quad k=1,2,\cdots,
\end{align*}
与假设矛盾。
\end{enumerate}

\end{proof}

\begin{theorem}[广义Minkowski不等式]\label{theorem:广义Minkowski不等式}
设 $f(x,y)$ 是 $\mathbb{R}^n\times\mathbb{R}^n$ 上的可测函数。若对几乎处处的 $y\in\mathbb{R}^n$，$f(x,y)\in L^p(\mathbb{R}^n)$（$1\leqslant p<+\infty$），且有
\begin{align*}
\int_{\mathbb{R}^n} \left( \int_{\mathbb{R}^n} |f(x,y)|^p \mathrm{d}x \right)^{\frac{1}{p}} \mathrm{d}y = M < +\infty,
\end{align*}
则
\begin{align}
\left( \int_{\mathbb{R}^n} \left| \int_{\mathbb{R}^n} f(x,y)\mathrm{d}y \right|^p \mathrm{d}x \right)^{\frac{1}{p}} \leqslant \int_{\mathbb{R}^n} \left( \int_{\mathbb{R}^n} |f(x,y)|^p \mathrm{d}x \right)^{\frac{1}{p}} \mathrm{d}y. \label{eq::8u238rURU32ur8-10}
\end{align}
\end{theorem}
\begin{proof}
不妨设 $p>1$，$p'$ 是 $p$ 的共轭指标，令
\begin{align*}
F(x) = \int_{\mathbb{R}^n} |f(x,y)|\mathrm{d}y.
\end{align*}
对任意可积简单函数 $\varphi(x)$，有
\begin{align*}
\left| \int_{\mathbb{R}^n} F(x)\varphi(x)\mathrm{d}x \right| &\leqslant \int_{\mathbb{R}^n} |F(x)\varphi(x)|\mathrm{d}x 
= \int_{\mathbb{R}^n} \left( \int_{\mathbb{R}^n} |f(x,y)|\mathrm{d}y \right) |\varphi(x)|\mathrm{d}x \\
&= \int_{\mathbb{R}^n} \left( \int_{\mathbb{R}^n} |f(x,y)| \cdot |\varphi(x)|\mathrm{d}x \right) \mathrm{d}y 
\leqslant \int_{\mathbb{R}^n} \left( \int_{\mathbb{R}^n} |f(x,y)|^p \mathrm{d}x \right)^{\frac{1}{p}} \left( \int_{\mathbb{R}^n} |\varphi(x)|^{p'} \mathrm{d}x \right)^{\frac{1}{p}'} \mathrm{d}y = M \|\varphi\|_{p'}.
\end{align*}
由\refthe{theorem:实变函数--定理6.19}，得 $\|F\|_p \leqslant M$，故
\begin{align*}
\left( \int_{\mathbb{R}^n} \left| \int_{\mathbb{R}^n} f(x,y)\mathrm{d}y \right|^p \mathrm{d}x \right)^{\frac{1}{p}} \leqslant \int_{\mathbb{R}^n} \left( \int_{\mathbb{R}^n} |f(x,y)|^p \mathrm{d}x \right)^{\frac{1}{p}} \mathrm{d}y.
\end{align*}

\end{proof}

\begin{remark}
在 $\mathbb{R}\times[0,2]$ 上定义二元函数
\begin{align*}
f(x,y) =
\begin{cases}
f(x), & 0\leqslant y<1, \\
g(x), & 1\leqslant y\leqslant 2,
\end{cases}
\end{align*}
其中 $f\in L^p(\mathbb{R})$，$g\in L^p(\mathbb{R})$，则广义Minkowski不等式化为通常的Minkowski不等式：
\begin{align*}
\|f+g\|_p \leqslant \|f\|_p + \|g\|_p.
\end{align*}
\end{remark}

\begin{remark}
在 $(0,+\infty)\times[1,2]$ 上定义二元函数
\begin{align*}
f(x,y) =
\begin{cases}
a_n, & n\leqslant x<n+1,\ 0\leqslant y<1, \\
b_n, & n\leqslant x<n+1,\ 1\leqslant y\leqslant 2,
\end{cases}
\quad (n\in\mathbb{N}),
\end{align*}
则广义Minkowski不等式化为
\begin{align*}
\left( \sum_{n=1}^\infty |a_n + b_n|^p \right)^{\frac{1}{p}} \leqslant \left( \sum_{n=1}^\infty |a_n|^p \right)^{\frac{1}{p}} + \left( \sum_{n=1}^\infty |b_n|^p \right)^{\frac{1}{p}}.
\end{align*}
\end{remark}

\begin{example}
设 $1<p<+\infty$，$f\in L^p((0,+\infty))$。若记
\begin{align*}
F(x) = \frac{1}{x} \int_0^x f(t)\mathrm{d}t, \quad x>0,
\end{align*}
则 $F\in L^p((0,+\infty))$，且 $\|F\|_p \leqslant \frac{p}{p-1}\|f\|_p$.
\end{example}

\begin{proof}
因为
\begin{align*}
F(x) = \int_0^1 f(xt)\mathrm{d}t,
\end{align*}
由\hyperref[theorem:广义Minkowski不等式]{广义Minkowski不等式}，得
\begin{align*}
\|F\|_p &= \left( \int_0^{+\infty} \left| \int_0^1 f(xt)\mathrm{d}t \right|^p \mathrm{d}x \right)^{\frac{1}{p}} \\
&\leqslant \int_0^1 \left( \int_0^{+\infty} |f(xt)|^p \mathrm{d}x \right)^{\frac{1}{p}} \mathrm{d}t = \int_0^1 \left( \int_0^{+\infty} |f(y)|^p t^{-1} \mathrm{d}y \right)^{\frac{1}{p}} \mathrm{d}t \\
&= \|f\|_p \int_0^1 t^{-\frac{1}{p}} \mathrm{d}t = \frac{p}{p-1}\|f\|_p.
\end{align*}
（当 $p=1$ 时，$F\notin L^1$。）

\end{proof}







































































\end{document}